% from ega4-21b-fr.tex
\oldpage[IV-4]{290}
\begin{env}[21.10.2]
\label{IV.21.10.2}
Nous nous proposons de montrer que:
\begin{enumerate}[label=\emph{\Alph*)}]
  \item Lorsque deux des conditions 1\textsuperscript{o}, 2\textsuperscript{o}, 3\textsuperscript{o} de \sref{IV.21.10.1} sont \emph{simultanément} satisfaites, les valeurs correspondantes de $n_{x'}$ \emph{coïncident}.
  \item L'ensemble des $x'\in X'^{(1)}$ tels que $n_{x'}\neq0$ est \emph{localement fini} dans $X'$.
\end{enumerate}

Pour démontrer A), supposons d'abord que $x\notin T$ et que $x$ vérifie l'une des conditions 2\textsuperscript{o} ou 3\textsuperscript{o} de \sref{IV.21.10.1};
alors $n_x=0$ et si l'on est dans le cas 2\textsuperscript{o}, on a $n_{x'}=0$;
si on est dans le cas 3\textsuperscript{o}, on a $Z_x=0$ puisque $\Supp(Z)\subset T$, donc encore $n_{x'}=0$.
Il reste à considérer le cas où l'on est à la fois dans le cas 2\textsuperscript{o} et dans le cas 3\textsuperscript{o};
alors, puisque $\dim(\sh{O}_{X,x})=1$, $\sh{O}_{X,x}$ est un anneau de valuation discrète;
si $t$ est une uniformisante de cet anneau, le diviseur correspondant à $Z_x$ est $\operatorname{div}(t^{n_x})$ dans $\operatorname{Div}(\sh{O}_{X,x})$, et son image dans $\operatorname{Div}(\sh{O}_{X',x'})$ est le diviseur de $t'^{n_x}$, où $t'$ est l'élément (régulier) de $\sh{O}_{X',x'}$ image de $t$.
On peut évidemment se borner au cas où $n_x=1$, et alors la définition \sref{IV.21.6.5.1} montre que l'image de $Z_x$ par \sref{IV.21.10.1.2} est le nombre $\lambda_{x'}$, ce qui achève de prouver A).

Démontrons maintenant B).
Posons $T_0=\Supp(Z)\subset T$, $T'_0=f^{-1}(T_0)$;
il suffit de prouver que la relation $n_{x'}\neq0$ implique que $x'$ appartient à l'ensemble des points maximaux de $T'_0$, ce dernier étant localement fini dans $X'$.
Il est immédiat que l'on a nécessairement $x'\in T'_0$;
si $x'$ n'était pas maximal dans l'ensemble fermé $T'_0$, il existerait
